\subsection{方阵的行列式}

定义2. 设 I 是有限集，A 是交换环，X 是环 A 上的 (I, I) 型方阵（II，§10，第 7 号）。A-模 A \(^{I}\) 的自同态 u（它关于 A \(^{I}\) 的典范基的矩阵为 X）的行列式称为 X 的行列式，记作 det X。

若 \(X = (\xi_{ij})_{(i,j) \in I \times I}\) 且 \((e_i)_{i \in I}\) 是 \(A^I\) 的典范基，则自同态 \(u\) 由下式给出

\[
u (e _ {i}) = \sum_ {j \in J} \xi_ {j i} e _ {j}.
\]

当 \(\mathbf{I} = [1, n] \subset \mathbf{N}\) 时，若记 \(x_{i} = u(e_{i})\) （对每个 \(i \in \mathbf{I}\) ），则 \(X\) 的行列式由关系式

\[
x _ {1} \wedge x _ {2} \wedge \dots \wedge x _ {n} = (\det X) e _ {1} \wedge e _ {2} \wedge \dots \wedge e _ {n}\tag{6}
\]

换言之，det X 等于行列式 \(\det(x_{1}, x_{2}, \ldots, x_{n})\) （关于 \(A^{n}\) 的典范基者）。因此：

命题 4. 对于 \(n\) 个向量 \(x_1, \ldots, x_n\) （ \(A^n\) 的），设 \(X(x_1, \ldots, x_n)\) 表示 \(n\) 阶方阵，其第 \(i\) 列为 \(x_i\) （ \(1 \leqslant i \leqslant n\) ）。则映射

\[
(x _ {1}, \dots , x _ {n}) \mapsto \det (X (x _ {1}, \dots , x _ {n}))
\]

（从 \((\mathbf{A}^n)^n\) 到 \(\mathbf{A}\) ）是交错的且 \(n\) -线性的。

特别地，若矩阵有两列相等，则其行列式为零。若对矩阵的列进行置换，则新矩阵的行列式等于原矩阵的行列式乘以 \(\varepsilon_{\sigma}\) 。若将矩阵的某一列加上另一个不同指标的列的标量倍数，则新矩阵的行列式等于原矩阵的行列式。

更一般地，设 M 是有限维数为 n 的自由 A-模，并设 \((e_{i})_{i\in I}\) 是 M 的一组基；对于 M 的每个自同构 u，若 X 是 u 关于基 \((e_{i})\) 的矩阵，则

\[
\det (u) = \det (X)\tag{7}
\]

这由定义直接推出。

当 \(\mathbf{I} = [1, n]\) 时， \(X\) 的行列式也记作

\[
\det (\xi_ {i j}) _ {1 \leqslant i \leqslant n, 1 \leqslant j \leqslant n}
\]

或在不致引起混淆时简记为 \(\operatorname{det}(\xi_{ij})\) ，也可记作

\[
\left| \begin{array}{c c c c c} \xi_ {1 1} & \xi_ {1 2} & \ldots & \xi_ {1 n} \\ \xi_ {2 1} & \xi_ {2 2} & \ldots & \xi_ {2 n} \\ \cdot & \cdot & \cdot & \cdot & \cdot \\ \xi_ {n 1} & \xi_ {n 2} & \ldots & \xi_ {n n} \end{array} \right|
\]

当 \(X = 1\) 时，矩阵 \(X\) 称为幺模的。

例。(1) 空矩阵的行列式等于 1；一阶方阵的行列式等于该矩阵的唯一元素。对于二阶矩阵

\[
\left( \begin{array}{c c} \xi_ {1 1} & \xi_ {1 2} \\ \xi_ {2 1} & \xi_ {2 2} \end{array} \right)
\]

则，在上述记号下，

\[
x _ {1} \wedge x _ {2} = (\xi_ {1 1} e _ {1} + \xi_ {2 1} e _ {2}) \wedge (\xi_ {1 2} e _ {1} + \xi_ {2 2} e _ {2}) = \xi_ {1 1} \xi_ {2 2} e _ {1} ^ {\prime} \wedge e _ {2} + \xi_ {2 1} \xi_ {1 2} e _ {2} \wedge e _ {1}
\]

，因此

\[
\left| \begin{array}{c c} \xi_ {1 1} & \xi_ {1 2} \\ \xi_ {2 1} & \xi_ {2 2} \end{array} \right| = \xi_ {1 1} \xi_ {2 2} - \xi_ {1 2} \xi_ {2 1}.
\]

我们将第 1 号和第 2 号中的一些结果翻译成矩阵的语言：

命题 5. 若 X 和 Y 是交换环 A 上具有相同有限指标集的两个方阵，则

\[
\det (X Y) = (\det X) (\det Y).\tag{8}
\]

\(X\) 可逆的必要且充分条件是 \(\det X\) 为 \(\mathbf{A}\) 的可逆元，且此时

\[
\det (X ^ {- 1}) = (\det X) ^ {- 1}.\tag{9}
\]

这直接由第 1 号命题 1 和第 2 号定理 1 得出。

推论。两个相似的方阵具有相等的行列式。

若 \(P\) 为可逆方阵，则由 (8) 和 (9) 可得 \(\det(PXP^{-1}) = \det X\) 。

命题6. 对于有限阶方阵X的列向量线性无关，其充分必要条件为det X在A中不是零因子。

这由第 2 号命题 3 得出。

